\subsection{Background}

Recent advancements in Artificial Intelligence (AI), particularly in Text-to-Speech (TTS) and Voice Conversion (VC) technologies, have enabled the generation of highly realistic synthetic speech capable of closely imitating human voices. Modern deep learning architectures and neural vocoders can now generate speech that is often indistinguishable from genuine human speech to human listeners. While these technologies have beneficial applications in accessibility, virtual assistants, and media production, they also introduce serious security and privacy concerns.

Synthetic speech and voice cloning technologies are increasingly being exploited in social engineering attacks, financial fraud, identity impersonation, misinformation campaigns, and unauthorized voice-based authentication bypass attempts. Existing deepfake speech detection systems generally perform well under controlled offline environments; however, their performance degrades significantly in real-world communication scenarios involving VoIP compression, codec distortion, packet loss, transmission noise, and streaming latency constraints.

As online communication systems continue to expand through platforms such as internet calling, conferencing, and voice-based authentication systems, there is a growing need for robust and low-latency detection mechanisms capable of operating under degraded communication conditions in near real-time environments.